% Minimal two-page skeleton: one question per page.
% Build: tectonic -X compile skeleton.tex
\documentclass[11pt,a4paper]{ctexart}
% Reusable preamble for a one-answer-per-page Chinese lab report.
% Requires: ctex (ctexart document class in the main file).
\usepackage[a4paper,left=2.4cm,right=2.4cm,top=2.8cm,bottom=2.3cm,headheight=14pt]{geometry}
\usepackage{amsmath,amssymb,bm}
\usepackage{booktabs,tabularx,array}
\usepackage{graphicx}
\usepackage{xcolor}
\usepackage[most]{tcolorbox}
\usepackage{fancyhdr}
\usepackage{enumitem}
\usepackage{siunitx}
\usepackage{pgfplots}
\usepgfplotslibrary{groupplots}
\pgfplotsset{compat=1.18}
\usepackage{circuitikz}
\usepackage{float}
\usepackage{caption}
\usepackage[hidelinks]{hyperref}
\captionsetup{font=small,labelfont=bf,skip=4pt}
\sisetup{detect-all}

\ctexset{
  section={format=\large\bfseries, beforeskip=0pt, afterskip=8pt},
  subsection={format=\normalsize\bfseries, beforeskip=6pt, afterskip=4pt}
}

% Header: page number == question number (enforced by one page per question).
\pagestyle{fancy}
\fancyhf{}
\fancyhead[L]{\small 实验五\quad 单管交流放大电路}
\fancyhead[R]{\small 第\thepage 题}
\renewcommand{\headrulewidth}{0.4pt}

\setlength{\parskip}{3pt}
\setlength{\parindent}{2em}
\setlength{\abovedisplayskip}{4pt}
\setlength{\belowdisplayskip}{4pt}
\setlength{\abovedisplayshortskip}{2pt}
\setlength{\belowdisplayshortskip}{3pt}

\newtcolorbox{qbox}{
  colback=black!4, colframe=black!4, boxrule=0pt, arc=1.2mm,
  left=4mm, right=4mm, top=2.6mm, bottom=2.6mm, before skip=2pt, after skip=10pt
}

% \question{5}{旁路电容对输出波形的影响}
\newcommand{\question}[2]{%
  \section*{第#1题\quad #2}%
}


\begin{document}

%====================== Q1 ======================
\question{1}{静态工作点}
\begin{qbox}
已知 $V_{CC}=12\,\mathrm{V}$，$R_C=2.4\,\mathrm{k\Omega}$，$R_E=1\,\mathrm{k\Omega}$，
$V_C=7.77\,\mathrm{V}$，$V_E=1.76\,\mathrm{V}$，求 $I_C$、$U_{CE}$。
\end{qbox}

由测量值可得
\begin{equation}
I_C=\frac{V_{CC}-V_C}{R_C}
=\frac{12-7.77}{2.4\times10^{3}}=1.76\,\mathrm{mA},
\end{equation}
\begin{equation}
U_{CE}=V_C-V_E=7.77-1.76=6.01\,\mathrm{V}.
\end{equation}

与记录值 $6.00\,\mathrm{V}$ 相差 $0.01\,\mathrm{V}$，属于读数误差。

%====================== Q2 ======================
\newpage
\question{2}{分压式偏置电路}
\begin{qbox}
画出分压式偏置共射放大电路，并说明 $R_E$ 的作用。
\end{qbox}

\begin{figure}[H]
\centering
\begin{circuitikz}[scale=0.78,transform shape]
  \node[npn] (Q) at (6,4) {};
  \draw (Q.C) -- (6,5.2) to[R, l=$R_C$] (6,6.6) -- (6,7.1) node[vcc]{$+V_{CC}$};
  \draw (Q.E) -- (6,2.8) to[R, l=$R_E$] (6,1.4) -- (6,1.0) node[ground]{};
  \draw (6,2.8) -- (7.5,2.8) to[C, l=$C_E$] (7.5,1.4) -- (6,1.4);
  \draw (Q.B) -- (4.6,4);
  \draw (4.6,4) to[R, l_=$R_{B1}$] (4.6,5.2) to[vR, l=$R_P$] (4.6,6.6) -- (6,6.6);
  \draw (4.6,4) -- (4.6,2.7) to[R, l=$R_{B2}$] (4.6,1.4) -- (4.6,1.0) node[ground]{};
  \draw (2.6,4) to[C, l=$C_1$] (4.6,4) (2.2,4) node[left]{$u_i$};
  \draw (6,5.2) -- (7.8,5.2) to[C, l=$C_2$] (9.0,5.2) node[right]{$u_o$};
\end{circuitikz}
\caption{分压式偏置共射放大电路}
\end{figure}

$R_E$ 引入直流负反馈：温度升高使 $I_C$ 增大时，$V_E$ 升高、$U_{BE}$ 减小，
从而抑制 $I_C$ 的变化，稳定静态工作点。

\end{document}
